problem:  
Let \(G\) be a compact, connected Lie group acting isometrically with cohomogeneity one on a closed manifold  
\[
M \cong [0,1]\times_{K_0,K_1} G/H,
\]
and suppose the isotropy representation of \(G/H\) splits as
\[
\mathfrak{m} = \mathfrak{m}_1 \oplus \mathfrak{m}_2
\]
into two inequivalent irreducible \(H\)-modules.  
Then each \(G\)-invariant metric on \(G/H\) can be written  
\[
g_t = x_1(t)\, Q|_{\mathfrak{m}_1} + x_2(t)\, Q|_{\mathfrak{m}_2},
\]
for smooth functions \(x_1,x_2>0\) and a fixed \(G\)-invariant inner product \(Q\) on \(\mathfrak{m}\); the metric on \(M\) reads
\[
g = dt^2 + g_t .
\]

Define the mixed intrinsic–extrinsic curvature functional
\[
\mathcal{F}_\lambda(g) = \int_0^1 \! \operatorname{Vol}(G/H,g_t) \Big( |\mathrm{Ric}_{g_t}|^2 + 
   \lambda |A_t|^2 \Big)\,dt,
\]
with \(\lambda \in \mathbb{R}\) and \(\operatorname{Vol}(M,g)=1\),  
where \(A_t=\frac12 g_t^{-1}\dot g_t\) is the shape operator of the orbits.

Let \((x_{1*},x_{2*})\) correspond to a homogeneous Einstein metric \(g_*\) on \(M\)
(i.e. \(g_t = r(t)^2(x_{1*},x_{2*})\) with fixed ratio \(x_{1*}:x_{2*}\)).  

---

**Precise conjectural problem (Linear stability and bifurcation for \(\mathcal{F}_\lambda\)):**

1.  Derive the Euler–Lagrange ODE system satisfied by \((x_1(t),x_2(t))\) for critical points of \(\mathcal{F}_\lambda\) under the given symmetry and boundary constraints.

2.  Linearize this system at the homogeneous Einstein solution \((x_{1*},x_{2*})\).  
    Denote by \(L_\lambda\) the linearized operator acting on \((\xi_1,\xi_2)\), the variations of \((x_1,x_2)\).

3.  **Conjecture (Spectral bifurcation criterion):**  
    There exists a discrete set of parameters \( \{\lambda_j\} \subset \mathbb{R} \) such that the kernel of \(L_\lambda\) changes dimension precisely at \(\lambda = \lambda_j\);  
    for those values, one can expect nontrivial branches of \(G\)-invariant critical metrics where the orbit‑ratio function \(x_1(t)/x_2(t)\) is non‑constant.  
    Equivalently, the orbit geometry ceases to be a fixed homogeneous metric up to scale and instead *bifurcates dynamically* along \(t\).

4.  **Rigidity question:**  
    For actions where the isotropy decomposition has only one irreducible summand (isotropy‑irreducible orbits), is \(L_\lambda\) always positive and no bifurcation possible?  
    That is, does the one‑summand case exhibit analytical rigidity for all \(\lambda\)?

Concrete test cases include the standard cohomogeneity‑one actions:
\[
SO(k)\times SO(n{+}1{-}k)\!\curvearrowright\! S^n,\qquad SU(2)\times SU(2)\!\curvearrowright\! S^7,
\]
for which explicit boundary smoothness conditions and two‑summand isotropy decompositions are known.

---

Why is it a "good" problem:  
This version is sharply formulated as a study of the *linearized variational operator* derived from an explicit geometric functional, under precise symmetry assumptions that guarantee a two‑parameter family of invariant orbit geometries.  
It recasts the vague idea of "symmetry enhancement" into concrete spectral language: the existence of bifurcation parameters \(\lambda_j\) where the operator \(L_\lambda\) loses invertibility.  
Proving or disproving the existence of such \(\lambda_j\) would reveal whether coupling intrinsic and extrinsic curvatures can destabilize homogeneous metrics, a genuinely open and analyzable question linking differential geometry, representation theory, and nonlinear ODE analysis.  

The problem thus achieves the desired mathematical precision (well‑defined variables, functional, operators, boundary constraints) while preserving profound depth: it tests whether analytical stability of geometric variational functionals can fail in a controlled symmetry‑reduced setting—an entryway to understanding how curvature variation interacts with global symmetry, a compelling frontier in geometric analysis.